\chapter{Risk Models}\label{ch:rs:risk-models}

A market-neutral book is forecast at 3.3 per cent annual volatility and realises 9.2. Its risk model knew the market, ten industries, beta, size and value,
and the book was neutral to all of them; it did not know momentum, and the book's proprietary score, without anyone having checked, was half momentum. Over
eight years its standardised returns have a standard deviation of 2.86 instead of one. The same book under a model that includes momentum is forecast at
9.23 per cent and realises 9.20. A risk model is a research tool before it is a reporting one: it says which risks a book takes, and it can
only say so for the risks it knows. This chapter builds \texttt{firm.riskmodel}, a fundamental equity model in the style that commercial vendors describe
publicly, fits it and a statistical model to \texttt{firm.synthmkt}, and evaluates them the way they should be evaluated, with \omterm{def:rs:risk-models:bias}{bias statistics}.

\section{Three families of factor models}

A factor model (Book~4, chapter~22) writes each stock's return as exposures times \omterm{def:rs:risk-models:families}{factor returns} plus a \omterm{def:rs:risk-models:families}{specific return}. Models differ in what they observe.

\begin{definition}[Fundamental, statistical and hybrid factor models; factor exposure, factor return, specific return]\label{def:rs:risk-models:families}
In a \emph{fundamental factor model}\index{fundamental factor model} the \emph{factor exposures}\index{factor exposure} are observed characteristics of each
stock (industry membership, standardised style scores) and the \emph{factor returns}\index{factor return} are estimated each period by cross-sectional
regression. In a \emph{statistical factor model}\index{statistical factor model} both are estimated from the returns alone, by principal components or
factor analysis. A \emph{hybrid factor model}\index{hybrid factor model} adds statistical factors to a fundamental model to capture what its characteristics
miss. The \emph{specific return}\index{specific return} is what the factors leave: the regression residual. (A third family, macroeconomic models, observes
factor returns such as inflation or the term spread and estimates exposures by time-series regression.)
\end{definition}

\begin{definition}[Style factor, industry factor]\label{def:rs:risk-models:style}
An \emph{industry factor}\index{industry factor} has as exposure a stock's membership of an industry; a \emph{style factor}\index{style factor} has as exposure a
standardised characteristic (size, value, momentum, an estimated beta), cap-weighted to mean zero and scaled to unit cross-sectional standard deviation.
\end{definition}

Rosenberg showed in 1974 that the covariance of stock returns has components beyond the market and that their loadings follow observable characteristics:
the origin of the fundamental model. Connor compared the three families on US returns; Chen, Roll and Ross priced macroeconomic factors. The chapter's model
is fundamental, the family commercial vendors document: for their US model, MSCI describes a Country factor that makes the \omterm{def:rs:risk-models:style}{industry factors} dollar-neutral
(each industry measured net of the market), \omterm{def:rs:risk-models:specific}{specific risk} estimated from daily time series and shrunk within size groups, and a \omterm{def:rs:risk-models:bias}{volatility regime adjustment}.

\section{Estimating factor returns}

Each day $t$ the model regresses the cross-section of returns on the exposures known at the previous close:
\[
\min_{f}\ \sum_i w_i\,(r_{it} - x_{i,t-1}^\top f)^2 \quad\text{subject to}\quad \sum_k c_k f_{k} = 0 \text{ over the industries},
\]
with weights $w_i$ the square root of capitalisation (larger stocks have smaller specific variance) and $c_k$ the capitalisation share of industry $k$. The
country factor's exposure is one for every stock, so the industry dummies alone would be collinear with it; the constraint makes the country factor the
cap-weighted market and each \omterm{def:rs:risk-models:style}{industry factor} its return net of the market. \texttt{firm.riskmodel.cs\_regression} solves the constrained problem through
its Karush--Kuhn--Tucker system (\cref{lst:rs:risk-models:csreg}).

On \texttt{firm.synthmkt}'s ten years, 1\,000 listed names a day, the model has fifteen factors: country, ten industries, and four styles, each known at
the previous close: beta, estimated over the trailing year and re-estimated monthly; size, the simulation's latent characteristic, fixed at listing; value,
the log of book to price from the latest filed book value (\texttt{firm.synthmkt.point\_in\_time\_styles}); momentum, the 12--1 month return, reset every
month. The country factor's annual volatility is 17.8\%, the industries' between 8.0\% and 8.7\%, momentum's 7.2\%. The average daily cross-sectional $R^2$
is 0.19: most of a stock's daily variance is specific, which is why a book of a hundred names a side still has \omterm{def:rs:risk-models:specific}{specific risk} worth modelling.

Point in time matters here as everywhere. The simulation keeps each stock's style exposures up to date in place; a first version of this chapter's model read
them at the end of the sample and used them for all ten years. It fitted without complaint, and its momentum factor had a volatility of 3.3\% instead of 7.2\%:
exposures from the future describe a different, smoother factor.

\section{Factor covariance}

The factor covariance forecast is an exponentially weighted estimate of the \omterm{def:rs:risk-models:families}{factor returns}, with a half-life of 90 days (Book~4, chapter~18, for the
exponentially weighted volatility). For horizons longer than a day, Newey--West terms add the autocovariances. Two refinements address the estimate's two
known failures. Volatility is not stationary: a backward-looking estimate lags when volatility rises. And sampling error makes the estimated covariance too
confident about some directions, which an optimiser then finds (Menchero, Wang and Orr: optimised portfolios' risk is underpredicted; chapter~25).

\begin{definition}[Bias statistic, volatility regime adjustment]\label{def:rs:risk-models:bias}
The \emph{bias statistic}\index{bias statistic} of a portfolio's risk forecasts over $T$ periods is the standard deviation of its standardised returns
$r_t/\hat\sigma_{t-1}$; it is one for correct forecasts, within $1 \pm \sqrt{2/T}$ with probability about 95\% for normal returns. The cross-sectional bias of
day $t$ is $B_t^2 = \frac1K\sum_k f_{kt}^2/\hat\sigma_{k,t-1}^2$ over the factors. A \emph{volatility regime adjustment}\index{volatility regime adjustment}
multiplies the covariance forecast by an exponentially weighted average $\lambda_t^2$ of recent $B_t^2$, raising all factor volatilities together when the day's
\omterm{def:rs:risk-models:families}{factor returns} are large for their forecasts.
\end{definition}

The chapter's adjustment uses a half-life of 42 days (\cref{lst:rs:risk-models:vra}). The simulated market's volatility follows a GJR-GARCH process, so regimes
come and go. On the cap-weighted market portfolio the overall \omterm{def:rs:risk-models:bias}{bias statistic} is 1.00 without the adjustment and 0.99 with it; the difference is in the regimes.
One-year rolling \omterm{def:rs:risk-models:bias}{bias statistics} range from 0.50 to 1.55 without it and from 0.56 to 1.32 with it, and 66\% of the one-year windows fall outside their band
($\pm 0.089$) without it against 44\% with it (\cref{fig:rs:risk-models:rolling}). A year is still a long time for a GARCH market; a desk that trades on daily
risk wants the adjustment and a short half-life.

\section{Specific risk}

\begin{definition}[Specific risk]\label{def:rs:risk-models:specific}
A stock's \emph{specific risk}\index{specific risk} is the forecast standard deviation of its \omterm{def:rs:risk-models:families}{specific return}.
\end{definition}

Each stock's specific variance is forecast by an exponentially weighted average of its squared residuals (half-life 90 days). Such forecasts are noisy, and
noise sorts: the stocks forecast most volatile are, on average, those whose recent residuals were luckily large. Standardising each stock's next-day residual
by its forecast and grouping by the forecast's decile shows it: the \omterm{def:rs:risk-models:bias}{bias statistic} is 1.13 in the lowest decile and 0.93 in the highest. Shrinking each
forecast volatility 10\% of the way towards the mean of its size decile, as MSCI describes for its model, flattens the pattern to between 0.97 and 1.04
(\cref{fig:rs:risk-models:deciles}). The remaining level, slightly above one, is the simulation's Student-$t$ shocks and earnings-day jumps.

\begin{omfigure}
\begin{tikzpicture}
\begin{axis}[width=10cm,height=5.2cm,ybar,bar width=7pt,xlabel={decile of forecast specific volatility},ylabel={bias statistic},tick label style={font=\small},
  label style={font=\small},xmin=0.4,xmax=10.6,ymin=0.85,ymax=1.15,xtick={1,2,3,4,5,6,7,8,9,10},grid=major,grid style={gray!25},legend columns=2,
  legend style={font=\footnotesize,at={(0.5,-0.3)},anchor=north,draw=none,fill=none,/tikz/every even column/.append style={column sep=8pt}},
  table/col sep=comma]
\addplot[fill=gray!45,draw=gray] table[x=decile,y=raw]{figdata/research/24-risk-models/deciles.csv};
\addplot[fill=omDef!55,draw=omDef] table[x=decile,y=shrunk]{figdata/research/24-risk-models/deciles.csv};
\draw[black!60,densely dotted] (axis cs:0.4,1) -- (axis cs:10.6,1);
\legend{time-series forecast,shrunk 10\% to its size decile}
\end{axis}
\end{tikzpicture}

\omcaption{\omterm{def:rs:risk-models:bias}{Bias statistics} of next-day \omterm{def:rs:risk-models:families}{specific returns} standardised by their forecasts, by decile of the forecast, years 3 to 10. Data:
\texttt{rs\_riskmodel.specific\_deciles}.}\label{fig:rs:risk-models:deciles}
\end{omfigure}

\section{Evaluating a risk model}

A portfolio's forecast variance is $x^\top F x + \sum_i w_i^2 s_i^2$, with $x = X^\top w$ its \omterm{def:rs:risk-models:families}{factor exposures}, $F$ the factor covariance and $s_i^2$ the
specific variances; each factor's contribution $x_k (Fx)_k$ sums to the factor part (\cref{lst:rs:risk-models:risk}). A model is judged by how its forecasts of
many portfolios' risk compare with what happened, not by the fit of its regression.

The chapter's test portfolios are held from each close over the next day and rebuilt monthly, and are scored over years 3 to 10 (2\,016 days; the bias band
is $1 \pm 0.031$). Fifty random market-neutral books, a hundred names a side, have \omterm{def:rs:risk-models:bias}{bias statistics} between 0.96 and 1.04, averaging 1.00; 92\% fall in the
band (98\% without the regime adjustment, which does nothing for books whose risk is mostly specific, and a little harm). The book of the opening is built
from a score equal to 0.5 times the momentum score plus $\sqrt{0.75}$ times noise, long the top hundred and short the bottom hundred, then neutralised to
every factor of the model without momentum. On day 1\,500 the full model forecasts it at 6.88\% a year, 76\% of the variance from its momentum exposure of
1.03; over the eight years its exposure has a root mean square of 1.23.

\begin{center}
\small
\begin{tabular}{@{}lccc@{}}
\toprule
the exposed book, years 3 to 10 & forecast & realised & \omterm{def:rs:risk-models:bias}{bias statistic}\\
\midrule
fundamental model, all factors & 9.23\% & 9.20\% & 1.04\\
fundamental model without momentum & 3.26\% & 9.20\% & 2.86\\
statistical model, 15 principal components & 5.59\% & 9.20\% & 1.68\\
\bottomrule
\end{tabular}
\end{center}

The model without momentum sees only the book's \omterm{def:rs:risk-models:specific}{specific risk}; the momentum factor's 7.2\% volatility times the book's exposure is the missing
two thirds. The statistical model, fifteen principal components refitted monthly on the trailing year, finds part of it without being told (5.59\%, a \omterm{def:rs:risk-models:bias}{bias
statistic} of 1.68): momentum's returns are a direction of common variation, which principal components look for, but its loadings are reset every month
while a trailing year's components assume loadings that stay put, so the direction the year's data describe is an average of twelve different ones. A
statistical model is still worth running beside the fundamental one as an alarm: here it forecasts 5.59\% where the fundamental model says 3.26\%, and when a
book's risk under the two models diverges, the fundamental model is missing something the book holds. That is the case for hybrid models too.

\begin{omfigure}
\begin{tikzpicture}
\begin{axis}[name=a,width=6.8cm,height=5.4cm,xlabel={\small year},ylabel={\small one-year bias statistic},tick label style={font=\footnotesize},
  xmin=3,xmax=10,ymin=0.4,ymax=3.6,grid=major,grid style={gray!25},title={\small the exposed book},legend columns=1,
  legend style={font=\scriptsize,at={(0.5,-0.32)},anchor=north,draw=none,fill=none},table/col sep=comma]
\addplot[omDef,very thick] table[x=year,y=full]{figdata/research/24-risk-models/rolling.csv};
\addplot[omThm,thick] table[x=year,y=nomom]{figdata/research/24-risk-models/rolling.csv};
\addplot[black!60,thick,dashed] table[x=year,y=pca]{figdata/research/24-risk-models/rolling.csv};
\draw[black!50,densely dotted] (axis cs:3,1) -- (axis cs:10,1);
\legend{all factors,no momentum,statistical}
\end{axis}
\begin{axis}[at={(a.east)},anchor=west,xshift=1.2cm,width=6.8cm,height=5.4cm,xlabel={\small year},tick label style={font=\footnotesize},
  xmin=3,xmax=10,ymin=0.4,ymax=1.8,grid=major,grid style={gray!25},title={\small the market portfolio},legend columns=1,
  legend style={font=\scriptsize,at={(0.5,-0.32)},anchor=north,draw=none,fill=none},table/col sep=comma]
\addplot[black!60,thick] table[x=year,y=mkt_raw]{figdata/research/24-risk-models/rolling.csv};
\addplot[omDef,very thick] table[x=year,y=mkt_vra]{figdata/research/24-risk-models/rolling.csv};
\draw[black!50,densely dotted] (axis cs:3,1) -- (axis cs:10,1);
\legend{no adjustment,regime adjustment}
\end{axis}
\end{tikzpicture}

\omcaption{One-year rolling \omterm{def:rs:risk-models:bias}{bias statistics}. Left: the book exposed to momentum under the three models. Right: the cap-weighted market under the full model,
with and without the \omterm{def:rs:risk-models:bias}{volatility regime adjustment}. Data: \texttt{rs\_riskmodel}.}\label{fig:rs:risk-models:rolling}
\end{omfigure}

\omterm{def:rs:risk-models:bias}{Bias statistics} have their own uncertainty. Over sixty months the band is $1 \pm 0.18$; with fat tails fewer than 95\% of correct forecasts fall inside it
(MSCI reports 86\% at a kurtosis of 5 over 120 periods). A single portfolio's \omterm{def:rs:risk-models:bias}{bias statistic} is weak evidence; many portfolios, rolling windows and the
portfolios the firm actually trades are what a model is judged on.

\section{Tutorial: predicted three, realised nine}\label{tut:rs:risk-models:tut}

\begin{tutorial}
\textbf{Goal.} Fit a fundamental model with and without momentum and a statistical model to \texttt{firm.synthmkt}, forecast the risk of books out of sample,
and evaluate by \omterm{def:rs:risk-models:bias}{bias statistics}. \textbf{End state:} the table, \cref{fig:rs:risk-models:rolling,fig:rs:risk-models:deciles}.
\begin{tutsteps}
\item \textbf{The daily regression} with the industry constraint.
\omcode{code/firm/riskmodel/firm_riskmodel.py}{30}{47}{Constrained weighted cross-sectional regression.}\label{lst:rs:risk-models:csreg}
\item \textbf{The \omterm{def:rs:risk-models:bias}{volatility regime adjustment}.}
\omcode{code/firm/riskmodel/firm_riskmodel.py}{87}{98}{The multiplier from the cross-sectional bias statistic.}\label{lst:rs:risk-models:vra}
\item \textbf{Portfolio risk} and its decomposition.
\omcode{code/firm/riskmodel/firm_riskmodel.py}{128}{134}{Factor and specific variance, and each factor's contribution.}\label{lst:rs:risk-models:risk}
\item \textbf{Run} \texttt{rs\_riskmodel.fundamental()}, \texttt{fundamental} with momentum omitted, \texttt{statistical()}, \texttt{summary()}, \texttt{random\_bias()},
\texttt{market\_bias()}, \texttt{specific\_deciles()} and \texttt{fig\_riskmodel.py}.
\end{tutsteps}
\textbf{What to change next.} Drop the value factor instead of momentum and build a book exposed to it; add three principal components of the fundamental
model's residuals as statistical factors (a hybrid model) and see whether the missing factor comes back.
\end{tutorial}

\section{Build: an equity risk model}\label{bld:rs:risk-models:riskmodel}

\begin{build}
\bfield{Purpose} The firm's equity risk model: exposures, \omterm{def:rs:risk-models:families}{factor returns}, covariance and \omterm{def:rs:risk-models:specific}{specific risk} every day, the risk of any book and its decomposition,
and the evidence (\omterm{def:rs:risk-models:bias}{bias statistics}) that the forecasts are right.
\bfield{Interface} \texttt{cs\_regression(r, X, w, C)}, \texttt{factor\_returns(R, exposures, W, C)}, \texttt{ewma\_cov(F, half\_life, nw\_lags)},
\texttt{vra(F, covs, half\_life)}, \texttt{specific\_var(E, half\_life, groups, shrink)}, \texttt{portfolio\_risk(w, X, Fcov, spec)}, \texttt{pca\_model(R, k)},
\texttt{bias\_stat(r, sigma)}, \texttt{bias\_band(T)}.
\bfield{Rules} Exposures known at the previous close; forecasts made at the close for the next period; every model change is judged on the \omterm{def:rs:risk-models:bias}{bias statistics} of
random books, of the factor portfolios and of the books the firm holds; a statistical model runs beside the fundamental one.
\bfield{Acceptance tests} \texttt{code/firm/riskmodel/tests/}: a constrained regression that recovers planted \omterm{def:rs:risk-models:families}{factor returns} and satisfies its constraint; the
covariance and the regime multiplier catching a planted volatility shift; specific variances recovering a planted level, shrinkage narrowing their dispersion;
contributions summing to the factor variance; orthonormal principal components; \omterm{def:rs:risk-models:bias}{bias statistics} of known forecasts and the band.
\bfield{Stretch} The eigenfactor adjustment of optimised portfolios' risk; hybrid statistical factors; multi-day horizons with Newey--West terms.
\end{build}

\begin{omsources}
\item B.~Rosenberg, ``Extra-market components of covariance in security returns'', \emph{Journal of Financial and Quantitative Analysis} 9(2), 1974.
\item G.~Connor, ``The three types of factor models: a comparison of their explanatory power'', \emph{Financial Analysts Journal} 51(3), 1995.
\item N.-F.~Chen, R.~Roll and S.~A.~Ross, ``Economic forces and the stock market'', \emph{Journal of Business} 59(3), 1986.
\item Y.~Liu, J.~Menchero, D.~J.~Orr and J.~Wang, \emph{The Barra US Equity Model (USE4): Empirical Notes}, MSCI, 2011.
\item J.~Menchero, J.~Wang and D.~J.~Orr, ``Improving risk forecasts for optimized portfolios'', \emph{Financial Analysts Journal} 68(3), 2012.
\end{omsources}

\section{Exercises}

\begin{exercise}[$\star$]\label{exo:rs:risk-models:1}
A book's forecast volatility is 3.26\% from \omterm{def:rs:risk-models:specific}{specific risk} only. Its exposure to an omitted factor of 7.2\% volatility has a root mean square of 1.23. What
should its realised volatility be, and what \omterm{def:rs:risk-models:bias}{bias statistic} follows?
\end{exercise}

\begin{exercise}[$\star$]\label{exo:rs:risk-models:2}
Give the 95\% band of the \omterm{def:rs:risk-models:bias}{bias statistic} for 60 monthly observations and for 2\,016 daily ones.
\end{exercise}

\begin{exercise}[$\star$]\label{exo:rs:risk-models:3}
Why does a model with a country factor and a full set of industry dummies need a constraint, and what does the cap-weighted constraint make the factors mean?
\end{exercise}

\begin{exercise}[$\star\star$]\label{exo:rs:risk-models:4}
Explain why time-series forecasts of specific volatility overpredict the high-volatility stocks and underpredict the low ones, even when each stock's volatility
is constant.
\end{exercise}

\begin{exercise}[$\star\star$]\label{exo:rs:risk-models:5}
The regime adjustment improved the market portfolio's rolling \omterm{def:rs:risk-models:bias}{bias statistics} and slightly worsened the random books'. Why?
\end{exercise}

\begin{exercise}[$\star\star$]\label{exo:rs:risk-models:6}
Show that the factor contributions $x_k(Fx)_k$ sum to the factor variance. Can a contribution be negative?
\end{exercise}

\begin{exercise}[$\star\star\star$]\label{exo:rs:risk-models:7}
\emph{Coding.} Build the exposed book from a score loading 0.3 instead of 0.5 on momentum, and report the three models' \omterm{def:rs:risk-models:bias}{bias statistics}.
\end{exercise}

\begin{exercise}[$\star\star\star$]\label{exo:rs:risk-models:8}
\emph{Find the flaw.} ``Our risk model explains only 19\% of daily return variance, so it is useless for a book of two hundred names.''
\end{exercise}

\section{Problem: Predicted Three, Realised Nine}

\begin{problem}[{Weekend problem --- a model that did not know}]\label{pb:rs:risk-models:1}
The chapter's models on \texttt{firm.synthmkt}.

\medskip\noindent\textbf{Part I --- The model.}
\begin{enumerate}
\item List the fifteen factors and how each exposure is obtained.
\item Why square-root-of-capitalisation weights?
\item What are the factor volatilities of the country, the industries and momentum?
\item What is the average cross-sectional $R^2$, and what does it mean for a book's risk?
\end{enumerate}

\medskip\noindent\textbf{Part II --- Covariance and \omterm{def:rs:risk-models:specific}{specific risk}.}
\begin{enumerate}[resume]
\item How does the \omterm{def:rs:risk-models:bias}{volatility regime adjustment} work, and what does it do to the market portfolio's rolling \omterm{def:rs:risk-models:bias}{bias statistics}?
\item Why does it not help the random books?
\item What pattern do raw specific forecasts show by decile, and what does shrinkage do?
\item Why does a residual level slightly above one remain?
\end{enumerate}

\medskip\noindent\textbf{Part III --- The exposed book.}
\begin{enumerate}[resume]
\item How is the exposed book built?
\item What are its forecasts, realised volatility and \omterm{def:rs:risk-models:bias}{bias statistics} under the three models?
\item What share of its forecast variance comes from momentum, at what exposure?
\item Why does the statistical model find most of the missing risk, but not all?
\end{enumerate}

\medskip\noindent\textbf{Part IV --- The verdict.}
\begin{enumerate}[resume]
\item State the \emph{named result}: the \omterm{def:rs:risk-models:bias}{bias statistic} of the model missing the planted factor, and the realised-to-predicted volatility ratio of the exposed
book.
\item How would the desk have found out before the drawdown?
\item What evidence would you require before adopting a new risk model?
\item How do fat tails change the reading of a \omterm{def:rs:risk-models:bias}{bias statistic}?
\item What is a \omterm{def:rs:risk-models:families}{hybrid factor model}, and when is it worth building?
\item Why are the random books' \omterm{def:rs:risk-models:bias}{bias statistics} a necessary but insufficient test?
\item Which book should a risk manager test first?
\item In one sentence: what does a risk model know?
\end{enumerate}
\end{problem}

\section{Interview questions}

\begin{interviewq}[$\star$ \iqroles{risk, researcher}]\label{iq:rs:risk-models:1}
What is the difference between a fundamental and a \omterm{def:rs:risk-models:families}{statistical factor model}? When would you prefer each?
\end{interviewq}

\begin{interviewq}[$\star\star$ \iqroles{researcher}]\label{iq:rs:risk-models:2}
How are the \omterm{def:rs:risk-models:families}{factor returns} of a fundamental model estimated each day, and why the constraint on the industries?
\end{interviewq}

\begin{interviewq}[$\star\star$ \iqroles{risk}]\label{iq:rs:risk-models:3}
A market-neutral book realises twice its forecast volatility. How do you find out why?
\end{interviewq}

\begin{interviewq}[$\star\star$ \iqroles{risk, researcher}]\label{iq:rs:risk-models:4}
How do you evaluate a risk model? What is a \omterm{def:rs:risk-models:bias}{bias statistic}, and what are its pitfalls?
\end{interviewq}

\begin{interviewq}[$\star\star$ \iqroles{researcher}]\label{iq:rs:risk-models:5}
Why do optimised portfolios tend to have their risk underpredicted?
\end{interviewq}

\begin{interviewq}[$\star\star\star$ \iqroles{researcher, developer}]\label{iq:rs:risk-models:6}
Design the daily production of an equity risk model: inputs, order of computation, checks, and what is stored.
\end{interviewq}
